\scanpage{27}
Or pour $u$ variable, $\Delta(u)$ est un polynôme
dans $\unsure{\mathbb Q_p|\mathbb Z_p|\deleted{\mathbb Q}\mathbb Z_p}[T]$ ($T=(T_1\mathbin{-}\mathbin{-}T_{\unsure{N|\alpha}})$, $N=\deleted{\unsure{\dim|\operatorname{rang}}}\operatorname{rg}_{\mathbb Z}\operatorname{End}(A)_{\mathrm{herm}}$)
qui, on vient de le voir, prend des valeurs entières
pour des valeurs entières des variables. Donc
$\Delta(u)$ est un polynôme à coeff. \deleted{\unsure{entiers}} dans $\mathbb Q$
(et non seulemt dans $\mathbb Q_p$), // \deleted{\unsure{donc même à coeff}
dans $\mathbb Z_{(p\mathbb Z)}$ [Car $\mathbb Z_p[T]\cap\mathbb Q[T]=\mathbb Z[T]$,
car $\mathbb Z_p\cap\mathbb Q=\mathbb Z_{(p\mathbb Z)}$]}. Il faut prouver que
ses valeurs sur le cône des $u>0$
\deleted{sont} $\geq0$, et \deleted{pour que} il suffit de le prouver
quand on passe de $\mathbb Q$ à $\mathbb R$. Mais alors
tt \unsure{élmt|élt} $\geq0$ est un carré, d'un élmt
hermitien \inserted{$v$} $>0$, donc $\Delta(u)=\Delta(v^2)=\Delta(v)^2\in\mathbb R^{2}=\mathbb R^+$,
cqfd.

b) $\alpha_p(\varphi)$ induit bien un \deleted{\unsure{hom}}\inserted{isom} de $\operatorname{d\acute et}T_p(A)_{\mathrm{\acute et}}$ avec
$\operatorname{d\acute et}T_p(A^*)_{\mathrm{\acute et}}$. Cela \deleted{\unsure{va déjà pr} \unsure{ces deux|les deux}
(i) $\alpha_p(\varphi)$ applique $\det T_p(A)_{\mathrm{\acute et}}$ sur $\det T_p(A^*)_{\mathrm{\acute et}}$ ;
i.e. $\beta_p(\varphi)$}\inserted{équivaut à la relation}
\[
v_p\bigl(\operatorname{card}(\operatorname{Coker}\deleted{\unsure{\beta_p}}T_p(\widetilde\varphi)_{\mathrm{\acute et}})\bigr)
=v_p\bigl(\sqrt{\deg\widetilde\varphi}\bigr)
\]
Or on a \unsure{évidemment|visiblement}
% archive sheet number: 26
